\documentclass[supercite]{Hust-undergraduate-kaiti}

\title{采用元数据键值对存储的文件操作设计与实现}
\school{计算机科学与技术学院}
\author{崔顺杰}
\classnum{CS2206}
\stunum{U202215522}
\instructor{曹强} % 李平、孙伟平、范晔斌、陈加忠
\usepackage{zhnumber} % change section number to chinese
\usepackage{tikz}
\usetikzlibrary{mindmap,trees}
\usepackage{algorithm, multirow}
\usepackage{algpseudocode}
\usepackage{amsmath}
\usepackage{amsthm}
\usepackage{framed}
\usepackage{mathtools}
\usepackage{subcaption}
\usepackage{amsmath, amssymb, amsfonts}
\usepackage{tabu}
\usepackage{xltxtra} %提供了针对XeTeX的改进并且加入了XeTeX的LOGO, 自动调用xunicode宏包(提供Unicode字符宏)
\usepackage{bm}
\usepackage{longtable}
\usepackage{tikz}
\usepackage{pdfpages}
\usepackage{tikzscale}
\usepackage{pgfplots}
%\usepackage{enumerate}
\usepackage{graphicx}
\pgfplotsset{compat=1.16}

\setCJKmainfont[
    BoldFont={FandolHei},
    ItalicFont={FandolKai}
]{FandolSong}
\setmainfont{Times New Roman}

\newcommand{\cfig}[3]{
  \begin{figure}[htb]
    \centering
    \includegraphics[width=#2\textwidth]{images/#1.tikz}
    \caption{#3}
    \label{fig:#1}
  \end{figure}
}

\newcommand{\sfig}[3]{
  \begin{subfigure}[b]{#2\textwidth}
    \includegraphics[width=\textwidth]{images/#1.tikz}
    \caption{#3}
    \label{fig:#1}
  \end{subfigure}
}

\newcommand{\xfig}[3]{
  \begin{figure}[htb]
    \centering
    #3
    \caption{#2}
    \label{fig:#1}
  \end{figure}
}

\newcommand{\rfig}[1]{\autoref{fig:#1}}
\newcommand{\ralg}[1]{\autoref{alg:#1}}
\newcommand{\rthm}[1]{\autoref{thm:#1}}
\newcommand{\rlem}[1]{\autoref{lem:#1}}
\newcommand{\reqn}[1]{\autoref{eqn:#1}}
\newcommand{\rtbl}[1]{\autoref{tbl:#1}}

\algnewcommand\Null{\textsc{null }}
\algnewcommand\algorithmicinput{\textbf{Input:}}
\algnewcommand\Input{\item[\algorithmicinput]}
\algnewcommand\algorithmicoutput{\textbf{Output:}}
\algnewcommand\Output{\item[\algorithmicoutput]}
\algnewcommand\algorithmicbreak{\textbf{break}}
\algnewcommand\Break{\algorithmicbreak}
\algnewcommand\algorithmiccontinue{\textbf{continue}}
\algnewcommand\Continue{\algorithmiccontinue}
\algnewcommand{\LeftCom}[1]{\State $\triangleright$ #1}


\newtheorem{thm}{定理}[section]
\newtheorem{lem}{引理}[section]
\renewcommand {\thefigure} {\arabic{figure}}
\colorlet{shadecolor}{black!15}

\theoremstyle{definition}
\newtheorem{alg}{算法}[section]
\def\thmautorefname~#1\null{定理~#1~\null}
\def\lemautorefname~#1\null{引理~#1~\null}
\def\algautorefname~#1\null{算法~#1~\null}
\usepackage{xcolor}
\usepackage{listings}
\usepackage[skins]{tcolorbox}
\tcbuselibrary{breakable}
\newtcolorbox{abox}[1][]{enhanced,
	fonttitle = \heiti \large \bfseries,% 标题字体设置
	colbacktitle =black,% 标题背景颜色
	coltitle=white,% 标题字体颜色
	% halign title=center,% 标题对齐方式
	attach boxed title to top left = {yshift = -5pt},
	%	将标题以box 形式放在文本框左上角，并向下移动5pt	
	%----------正文字体--------
	fontupper = \kaishu,%正文字体设置
	%colback = white,%正文背景颜色设置
	%---------边框设置---------
	arc = 4pt,%正文边框边角弧度
	%	boxed title style={arc=0pt},%标题框边角弧度
	colframe=black,%边框颜色设置
	toprule  = 1pt,%取消上边框
	rightrule = 1pt,%取消下边框
	%---------调整字体位置------
	top = 10pt,%增大正文文本与上边框距离
	%---------设置标题为默认参数，不影响其他参数的设置
	title=#1,
	breakable
}
\newcommand{\tabincell}[2]{\begin{tabular}{@{}#1@{}}#2\end{tabular}}  
\begin{document}

\maketitle
\clearpage

\pagenumbering{Roman}

\pagenumbering{arabic}
\section{课题来源、目的及意义}
\subsection{课题来源}

本课题来源于国家自然科学基金项目"面向新型混合存储的高效异构融合系统架构及机制"（62172175）以及国家重点研发计划"ZB级海量冷数据存储架构及高效存储管理系统"（2024YFB4505105）。

随着云计算、大数据和人工智能应用的快速发展，现代存储系统面临前所未有的元数据管理压力。据Meta公司2020年公开的数据中心工作负载分析，元数据操作占据了超过50\%的总I/O操作量；在海量小文件场景下（如AI训练数据集、容器镜像、日志采集等），元数据开销更是高达总延迟的40\%\cite{ref_bfo}。传统文件系统（ext4、XFS等）的元数据管理方式已难以满足高并发、低延迟的性能需求。近年来，JuiceFS、CephFS、Hadoop Ozone等工业级系统纷纷采用键值（Key-Value, KV）存储引擎管理元数据，证明了这一技术路线的可行性。本课题旨在设计并实现文件元数据的键值对存储方案和相应的POSIX文件操作函数。

\subsection{课题目的}

本课题的核心目的是：设计并实现文件元数据的键值对存储和相应的主要POSIX文件操作函数。具体而言包含以下三个层面：

（1）设计文件命名空间到KV存储的映射方案。将文件系统的层级命名空间（目录树、inode属性、目录条目等）高效地映射到KV存储的扁平键空间上。

（2）实现标准POSIX文件操作语义。在KV存储映射基础上，实现lookup、getattr、setattr、readdir、create、mkdir、unlink、rmdir、rename、open、read、write等标准文件操作，支持通过Linux FUSE进行挂载和透明访问。

（3）构建可运行的原型系统。

\subsection{课题意义}

\subsubsection{理论意义}

传统文件系统使用树形结构管理元数据，在大量小文件创建、删除等场景下存在严重的随机I/O瓶颈。以ext4为例，一次create操作需修改inode位图、inode表、目录HTree、Journal等4$\sim$6个分散的磁盘区域，全部为随机写。而LSM-tree（Log-Structured Merge Tree）作为KV存储的核心数据结构，将随机写入转换为顺序追加写入，理论上可大幅提升元数据写入性能\cite{ref_lsm}。本课题通过实际系统验证这一理论优势，为文件系统元数据管理提供新的技术路径。

\subsubsection{实践意义}

当前工业界对基于KV存储的文件元数据管理有强烈需求，但这些系统设计复杂度极高，缺乏对"KV存储如何承载POSIX语义"这一核心问题的系统性分析。本课题通过从零构建一个完整原型系统，为理解和优化KV元数据管理提供参考实现。

\section{国内外研究现况及发展趋势}
\subsection{基于KV存储的文件元数据管理}

利用KV存储管理文件系统元数据的研究始于2013年卡内基梅隆大学的TableFS系统，此后十余年间形成了从学术原型到超大规模生产系统的完整技术演进链。

\subsubsection{学术验证阶段（2013—2017）}

TableFS\cite{ref_tablefs}首次将LevelDB嵌入本地文件系统管理全量元数据，100万文件创建比ext4快5倍，证明了KV存储替代传统文件系统管理元数据的可行性。IndexFS\cite{ref_indexfs}将这一理念扩展到分布式场景，通过GIGA+哈希分裂算法支持128个MDS节点的水平扩展。LocoFS\cite{ref_locofs}通过目录元数据与文件属性的松耦合解聚合，将KV引擎利用率从18\%提升到93\%。

\subsubsection{工业落地阶段（2018—2020）}

JuiceFS使用Redis/TiKV管理元数据，在生产环境中广泛验证了KV元数据管理的可行性。Hadoop Ozone用三级RocksDB突破了HDFS 10亿文件限制。Yang等人的BFO（Batch-File Operations）\cite{ref_bfo}研究深入分析了批量文件操作中的元数据开销，发现80\%的小文件操作目标小于32字节，元数据开销占总延迟的40\%，这一发现直接启发了本课题中WriteBatch批量提交机制的设计。

\subsubsection{极致优化阶段（2023—2025）}

SingularFS\cite{ref_singularfs}证明精心优化的单机元数据服务器性能可超过多节点分布式方案，单MDS达8.36M IOPS，打破了"必须分布式"的思维定势——这对本课题的启示是，在走向分布式前应充分挖掘单机性能潜力。Mantle\cite{ref_mantle}提出Delta Record追加更新与IndexNode内存缓存层，避免高频属性的读-改-写模式，吞吐提升最高115倍。本课题的Delta增量更新机制直接受Mantle启发。

\subsection{发展趋势}

综合以上分析，文件元数据KV存储技术呈现以下趋势：

（1）事务化与原子性保证成为标配，从简单KV读写演进到WriteBatch局部事务保证；

（2）元数据与数据分离架构成为主流，实现独立优化和扩展；

（3）增量更新替代读-改-写，通过追加式Delta记录减少写放大；

（4）单机性能挖掘仍有巨大空间，SingularFS等系统证明单机可达十亿级文件管理。

\section{研究内容与技术方案}

\subsection{预计达到的目标}

能够通过标准POSIX文件接口进行文件/目录访问，并进一步优化整体性能。具体包括：完成文件命名空间到KV存储的完整映射设计与实现；实现create、mkdir、lookup、readdir、rename、unlink、rmdir等22种FUSE文件操作；构建事务化元数据管理引擎，保证所有元数据写操作的原子性；在实际硬件平台上与ext4进行性能对比测试。

\subsection{系统整体架构}

\begin{figure}[htb]
    \centering
    \includegraphics[width=0.6\textwidth]{overall.png}
    \caption{RucksFS系统整体架构}
    \label{fig:overall}
\end{figure}

系统使用Rust语言开发，采用Cargo Workspace组织为多个独立crate。整体架构如图\ref{fig:overall}所示，为单机FUSE文件系统，自上而下分为四个层次：

（1）\textbf{FuseClient层。}实现fuser库\cite{ref_fuser}的Filesystem trait，负责解析Linux内核经/dev/fuse设备转发的VFS请求，将其转换为内部VfsOps trait调用。用户程序的POSIX syscall经内核VFS和FUSE内核模块透明转发至此层。

（2）\textbf{VfsCore层。}作为元数据操作与数据操作的路由中心，将VfsOps请求拆分为MetadataOps（元数据读写）和DataOps（文件内容I/O）两类调用，分别下发至对应的服务模块。

（3）\textbf{Server层。}包括MetadataServer和DataServer两个核心组件。MetadataServer负责事务化元数据操作、Delta增量更新、LRU缓存和后台Compaction；DataServer负责文件数据的读写。

（4）\textbf{Storage层。}MetadataServer对接RocksDB\cite{ref_rocksdb}存储引擎，使用四个Column Family分别管理不同类型的元数据；DataServer对接RawDisk进行文件内容存储。通过MetadataOps、DataOps、VfsOps三层Trait抽象，上层逻辑与底层存储后端解耦，后端可独立替换。

\subsection{元数据KV存储设计}

本课题的核心工作是设计文件命名空间到KV存储的映射方案。以下从映射策略选择、Column Family设计、键编码方案和Delta增量更新四个方面详细阐述。

\subsubsection{层级KV映射方案}

将文件系统的层级命名空间映射到KV存储的扁平键空间，存在两种主流方案。方案A为扁平KV映射，以文件完整路径（如\texttt{"/home/alice/notes.txt"}）作为Key，Value为文件属性。方案B为层级KV映射，以（父目录inode编号 + 子项名称）作为Key，仿照传统文件系统的目录项结构。

本系统采用方案B，原因在于三点。第一，rename目录时，扁平方案需要重写所有子孙路径前缀的Key，复杂度为$O(N)$（$N$为子孙文件数），而层级方案只需修改一条目录条目，复杂度为$O(1)$。第二，readdir时，层级方案以固定8字节的父目录inode编号为前缀做精确扫描，不会像扁平方案那样误扫到名称相似的兄弟目录（如\texttt{"/home/alice"}会误命中\texttt{"/home/alice\_backup/"}）。第三，层级方案的Key平均长度约20字节（8字节inode加文件名），远小于扁平方案30$\sim$100字节的平均路径长度，存储更紧凑。

\subsubsection{Column Family设计}

所有元数据存储在单个RocksDB实例中，使用四个Column Family（CF）进行逻辑隔离，各CF拥有独立的MemTable、SSTable和Compaction策略：

\begin{table}[htb]
\centering
\caption{RocksDB Column Family设计}
\label{tbl:cf}
\begin{tabular}{|c|p{2cm}|p{4.5cm}|p{4.5cm}|}
\hline
\textbf{CF名称} & \textbf{用途} & \textbf{Key格式} & \textbf{Value格式} \\
\hline
inodes & inode属性 & inode编号（8字节大端序） & 序列化的FileAttr结构体 \\
\hline
dir\_entries & 目录子项 & 父inode（8字节大端序）+ 子项名称（UTF-8） & 子inode编号（8字节）+ 文件类型（4字节） \\
\hline
delta\_entries & 属性增量 & inode（8字节大端序）+ 序列号（8字节大端序） & 序列化的DeltaOp（5$\sim$9字节） \\
\hline
system & 系统计数器 & ASCII字符串（如\texttt{"next\_inode"}） & 计数器值（8字节） \\
\hline
\end{tabular}
\end{table}

四个CF分离而非合并为一个的原因是：inodes CF以点查为主，适合配置Bloom Filter加速；dir\_entries CF以前缀扫描为主，适合配置Prefix Extractor；delta\_entries CF追加写密集，独立Compaction可避免干扰inodes CF中稳定的基础值；system CF极少访问，独立后不影响其他CF的性能。

\subsubsection{键编码规则}

所有多字节整数字段均采用大端序编码。这一设计保证了RocksDB默认的字典序字节比较能产生正确的数值排序，从而支持前缀扫描和范围查询。具体编码方式为：inodes CF的Key为8字节大端序的inode编号；dir\_entries CF的Key由8字节大端序的父inode编号紧接UTF-8编码的子项名称组成，无需分隔符（因前缀为固定8字节长度）；delta\_entries CF的Key由8字节大端序的inode编号加8字节大端序的全局递增序列号组成，序列号保证同一inode的多条Delta按写入顺序排列。

\subsubsection{Delta增量更新机制}

每次create、mkdir、unlink、rmdir、rename等操作都必须更新父目录inode的nlink（子项计数）、mtime（修改时间）和ctime（变更时间）三个属性。传统方式需要在事务中对父inode加行锁、读取57字节的完整属性、反序列化、修改、序列化、写回——同一目录下所有并发操作被串行化在这把锁上，成为热目录的首要瓶颈。

受Mantle\cite{ref_mantle}的Delta Record机制启发，本系统采用追加写方案：不修改父inode本体，而是向delta\_entries CF追加一条5$\sim$9字节的增量记录，描述"对父inode做何种变化"。追加操作无需读取父inode，也无需事务行锁，$N$个并发线程可以同时写入。

读取inode属性时采用Fold-on-Read策略：首先查询LRU缓存，命中则$O(1)$返回；未命中则读取inodes CF中的基础值，并扫描delta\_entries CF中该inode的所有增量记录，将基础值与所有Delta逐条折叠合并得到最终属性值，随后写入缓存。后台Compaction Worker每5秒扫描一次，当某inode的Delta累积超过32条时触发合并：将Delta折叠进基础值，原子写回inodes CF，清除旧Delta记录并使缓存失效。这一机制保证了读放大始终有界。

\subsection{事务化POSIX操作设计}

在KV映射方案的基础上，本系统使用RocksDB的PCC（Pessimistic Concurrency Control）事务模式实现POSIX文件操作的原子性和隔离性。每次写操作通过GetForUpdate获取行级排他锁，所有修改先缓存在内存WriteBatch中，commit时一次性写入WAL和MemTable，支持跨Column Family的原子写入。加锁失败立即返回Busy而非阻塞等待，上层封装execute\_with\_retry自动重试最多3次。

以三个核心操作为例说明事务设计模式：

（1）\textbf{create操作。}开启事务后，以GetForUpdate锁定目标目录条目Key，若已存在则返回EEXIST错误并回滚；否则分配新inode编号，向inodes CF写入新inode属性，向dir\_entries CF写入新目录条目，commit提交事务。事务提交后，在事务外向delta\_entries CF追加父目录的nlink+1和mtime/ctime更新记录。

（2）\textbf{rename操作（跨目录）。}开启事务后，对源目录条目和目标目录条目的Key按inode编号排序后依次GetForUpdate加锁（排序加锁防止死锁），验证源存在、处理目标已存在的覆盖情况，执行delete源条目 + put目标条目，commit后追加源父目录和目标父目录的Delta更新。

（3）\textbf{rmdir操作。}开启事务后，GetForUpdate锁定目录条目和目录inode，随后在事务内执行prefix\_scan检查目录是否为空（事务内扫描保证TOCTOU安全），非空则返回ENOTEMPTY，为空则delete目录条目和inode，commit后追加父目录Delta。

上述设计的关键在于：事务只保护结构变更（inode和dir\_entry的原子写入），父目录属性更新通过Delta在事务外追加完成，从而缩小事务临界区，提高同目录下不同文件操作的并行度。

\subsection{测试与评价方案}

测试分为两个层次：POSIX正确性验证和元数据性能评测。正确性方面，使用pjdfstest标准合规套件覆盖chmod、link、mkdir、rename、rmdir、unlink等操作的边界情况，配合Rust集成测试覆盖正常与异常路径。性能方面，使用自研Rust压测工具和mdtest标准工具，在10K/100K/1M文件梯度和1/2/4/8线程并发下与ext4进行对比评测。

\section{课题研究进度安排}

本课题的研究进度安排如表\ref{tbl:schedule}所示。

\begin{table}[htb]
\centering
\caption{课题研究进度安排}
\label{tbl:schedule}
\begin{tabular}{|c|c|p{9cm}|}
\hline
\textbf{阶段} & \textbf{时间} & \textbf{工作内容} \\
\hline
第一阶段 & 第7学期 11—12月 & 文献调研：研读TableFS、IndexFS、LocoFS、SingularFS、BFO、Mantle等论文；撰写技术调研报告；明确系统架构设计 \\
\hline
第二阶段 & 第7学期 1—2月 & 核心实现（一）：完成KV编码方案设计与实现、RocksDB存储后端、MetadataServer基础框架、基本POSIX操作（create/mkdir/lookup/getattr/readdir/unlink/rmdir） \\
\hline
第三阶段 & 第8学期 2—3月 & 核心实现（二）：完成Delta增量更新与后台Compaction、rename/setattr等复杂操作、PCC事务化元数据管理、FUSE客户端集成 \\
\hline
第四阶段 & 第8学期 3—4月 & 测试与优化：部署pjdfstest正确性验证、自研压测工具和mdtest性能评测、并发安全修复、WAL写入优化、锁粒度细化 \\
\hline
第五阶段 & 第8学期 4—5月 & 论文撰写：整理实验数据、完成与ext4的性能对比分析、撰写毕业论文、准备答辩材料 \\
\hline
\end{tabular}
\end{table}

\newpage

\section{主要参考文献}
\renewcommand{\refname}{}
\vspace{-2em}
\begin{thebibliography}{99}

\bibitem{ref_bfo}Yang Yang, Qiang Cao, Li Yang, et al. Batch-File Operations to Optimize Massive Files Accessing: Analysis, Design, and Application[J]. ACM Transactions on Storage, 2020, 16(3): 1-33.

\bibitem{ref_singularfs}Jing Liu, Andrea Arpaci-Dusseau, Remzi Arpaci-Dusseau, et al. SingularFS: A Billion-Scale Distributed File System Using a Single Metadata Server[C]. Proceedings of the 2023 USENIX Annual Technical Conference (ATC'23), 2023.

\bibitem{ref_fetchbpf}Xuechun Cao, Shaurya Patel, Soo-Yee Lim, et al. FetchBPF: Customizable Prefetching Policies in Linux with eBPF[C]. Proceedings of the 2024 USENIX Annual Technical Conference (ATC'24), 2024.

\bibitem{ref_fbmm}Sandeep Kumar, Aravinda Prasad, et al. FBMM: Making Memory Management Extensible With Filesystems[C]. Proceedings of the 2024 USENIX Annual Technical Conference (ATC'24), 2024.

\bibitem{ref_tablefs}Kai Ren, Garth Gibson. TableFS: Enhancing Metadata Efficiency in the Local File System[C]. Proceedings of the 2013 USENIX Annual Technical Conference (ATC'13), 2013.

\bibitem{ref_indexfs}Kai Ren, Qing Zheng, Swapnil Patil, et al. IndexFS: Scaling File System Metadata Performance with Stateless Caching and Bulk Insertion[C]. Proceedings of the IEEE/ACM International Conference on High Performance Computing, Networking, Storage and Analysis (SC'14), 2014.

\bibitem{ref_locofs}Siyang Li, Youyou Lu, Jiwu Shu, et al. LocoFS: A Loosely-Coupled Metadata Service for Distributed File Systems[C]. Proceedings of the IEEE/ACM International Conference on High Performance Computing, Networking, Storage and Analysis (SC'17), 2017.

\bibitem{ref_mantle}Wenhao Lv, Qiang Cao, et al. Mantle: A Scalable Hierarchical Namespace Service for Object Storage[C]. Proceedings of the ACM SIGOPS Symposium on Operating Systems Principles (SOSP'25), 2025.

\bibitem{ref_rocksdb}Facebook Inc. RocksDB: A Persistent Key-Value Store for Fast Storage Environments[EB/OL]. https://rocksdb.org, 2013.

\bibitem{ref_lsm}Patrick O'Neil, Edward Cheng, Dieter Gawlick, et al. The Log-Structured Merge-Tree (LSM-Tree)[J]. Acta Informatica, 1996, 33(4): 351-385.

\bibitem{ref_fuser}libfuse Contributors. fuser: Filesystem in Userspace (FUSE) for Rust[EB/OL]. https://github.com/cberner/fuser, 2020.

\end{thebibliography}
\clearpage
\begin{comment}
%\centerline{\zihao{4} \heiti 华中科技大学本科生毕业设计（论文）开题报告评审表}
%\thispagestyle{empty}
% Please add the following required packages to your document preamble:
% \usepackage{multirow}
% \usepackage[table,xcdraw]{xcolor}
% If you use beamer only pass "xcolor=table" option, i.e. \documentclass[xcolor=table]{beamer}
%
%\begin{table}[htbp]

%\resizebox{150mm}{93mm}{
%\setlength{\tabcolsep}{10mm}{

%\begin{tabular}{|cccccc|}
%\hline
%\multicolumn{1}{|c|}{\textbf{姓名}} & \multicolumn{1}{c|}{\textbf{}} & \multicolumn{1}{c|}{\textbf{学号}} & \multicolumn{1}{c|}{\textbf{}} & \multicolumn{1}{c|}{\textbf{指导教师}} & \textbf{} \\ \hline
%\multicolumn{2}{|c|}{\textbf{院（系）专业}}                              & \multicolumn{4}{c|}{\textbf{}}                                                                                     \\ \hline
%\multicolumn{6}{|c|}{\textbf{指导教师评语}}       \\                                                                            \multicolumn{6}{|l|}{{1.学生前期表现情况。}}    \\

%\multicolumn{6}{|l|}{{2.是否具备开始设计（论文）条件？是否同意开始设计（论文）？ }}   \\
%\multicolumn{6}{|l|}{{3.不足与建议。}}  \\
%\hline
%\multicolumn{6}{|c|}{{\color[HTML]{FF0000} }}                                                                                                                                           \\
%\multicolumn{6}{|c|}{{\color[HTML]{FF0000} }}                                                                                                                                           \\
\multicolumn{6}{|c|}{{\color[HTML]{FF0000} }}                                                                                                                                           \\
\multicolumn{6}{|c|}{{\color[HTML]{FF0000} }}                                                                                                                                           \\
\multicolumn{6}{|c|}{{\color[HTML]{FF0000} }}                                                                                                                                           \\
\multicolumn{6}{|c|}{{\color[HTML]{FF0000} }}                                                                                                                                                                                                                    \\
\multicolumn{6}{|c|}{{\color[HTML]{FF0000} }}                                                                                                                                           \\
\multicolumn{6}{|c|}{{\color[HTML]{FF0000} }}                                                                                                                                           \\
\multicolumn{6}{|c|}{\multirow{-15}{*}{{\color[HTML]{FF0000} }}}                                                                                                                        \\ \hline
\multicolumn{6}{|l|}{}                                                                                                                                                                  \\
\multicolumn{6}{|l|}{\multirow{-2}{*}{\qquad \qquad \qquad \qquad \qquad \qquad \qquad \qquad 指导教师（签名）：}}                                                                                                                                       \\ 
\multicolumn{6}{|r|}{年\qquad 月\qquad 日}                                                                                                                                                       \\ \hline
\multicolumn{6}{|c|}{\textbf{教研室（系、所）或开题报告答辩小组审核意见}}                                                                                                                                    \\ \hline
\multicolumn{6}{|c|}{{\color[HTML]{FF0000} }}                                                                                                                                           \\
\multicolumn{6}{|c|}{{\color[HTML]{FF0000} }}                                                                                                                                           \\
\multicolumn{6}{|c|}{{\color[HTML]{FF0000} }}                                                                                                                                           \\
\multicolumn{6}{|c|}{{\color[HTML]{FF0000} }}                                                                                                                                           \\
\multicolumn{6}{|c|}{\multirow{-5}{*}{{\color[HTML]{FF0000} }}}                                                                                                                         \\ \hline
\multicolumn{6}{|l|}{}                                                                                                                                                                  \\
\multicolumn{6}{|l|}{\multirow{-2}{*}{教研室（系、所）或开题报告答辩小组负责人（签名）：}}                                                                                                                       \\ 
\multicolumn{6}{|r|}{年\qquad 月\qquad 日}                                                                                                                                                       \\ \hline
\end{tabular}
}
}

\end{table}
\end{comment}
\includepdf[]{认定表.pdf}
\end{document}
